% Compléments Bloc 4 — zooms management d'équipe

\begin{frame}{C4.3 — De la capacité prévue à la capacité réelle}
\context{B4}{C4.3 · Lynx Immo · Zoom preuve}
\begin{block}{Écart observé}
Le planning supposait une certaine capacité de travail, mais la disponibilité réelle variait avec les cours, l'apprentissage et les contraintes individuelles.
\end{block}
\begin{tightitems}
  \item Rendre les tâches plus petites et plus explicites.
  \item Clarifier les dépendances pour éviter les blocages en chaîne.
  \item Réassigner, différer ou simplifier certains éléments.
  \item Intégrer le temps d'apprentissage dans la charge réelle du projet.
\end{tightitems}
\begin{block}{Message jury}
Coordonner l'équipe, c'est maintenir une trajectoire réaliste lorsque le plan initial rencontre les contraintes humaines du projet.
\end{block}
\end{frame}

\begin{frame}{C4.4 — Accompagner sans centraliser tout le travail}
\context{B4}{C4.4 · Lynx Immo · Zoom preuve}
\begin{columns}[T]
\begin{column}{0.48\linewidth}
\begin{block}{Mécanismes d'accompagnement}
\begin{tightitems}
  \item Explications techniques et prototypes.
  \item Documentation des décisions.
  \item Clarification des tâches et critères d'acceptation.
  \item Aide ciblée lorsqu'un membre reste bloqué.
\end{tightitems}
\end{block}
\end{column}
\begin{column}{0.48\linewidth}
\begin{block}{Limite à préserver}
\begin{tightitems}
  \item Aider n'est pas reprendre systématiquement la tâche.
  \item L'objectif est de rendre le membre suffisamment autonome pour livrer sa responsabilité.
  \item Il faut équilibrer aide, délégation et contraintes du projet.
\end{tightitems}
\end{block}
\end{column}
\end{columns}
\end{frame}

\begin{frame}{C4.6 — Évaluer avec des critères explicites, pas avec une impression}
\context{B4}{C4.6 · Lynx Immo · Zoom preuve}
\begin{columns}[T]
\begin{column}{0.47\linewidth}
\begin{block}{Collectif}
\begin{tightitems}
  \item Avancement par rapport à la roadmap.
  \item Couverture des exigences.
  \item Blocages et dépendances résolues.
  \item Qualité des livrables communs.
\end{tightitems}
\end{block}
\end{column}
\begin{column}{0.49\linewidth}
\begin{block}{Individuel dans le cadre projet}
\begin{tightitems}
  \item Responsabilités prises en charge.
  \item Capacité à progresser et communiquer.
  \item Contribution aux objectifs du projet.
  \item Feedback de fin de projet à double sens.
\end{tightitems}
\end{block}
\end{column}
\end{columns}
\begin{block}{Nuance}
Il s'agit d'un suivi de contribution dans une équipe étudiante, pas d'une évaluation RH d'un salarié.
\end{block}
\end{frame}

\begin{frame}{Bloc 4 — Ce que je transposerais dans un contexte professionnel complet}
\context{B4}{Projection professionnelle}
\begin{columns}[T]
\begin{column}{0.48\linewidth}
\begin{block}{Ce que je conserverais}
\begin{tightitems}
  \item Compétences dérivées des besoins produit.
  \item Responsabilités explicites.
  \item Feedback régulier et bidirectionnel.
  \item Montée en compétence incluse dans la charge.
\end{tightitems}
\end{block}
\end{column}
\begin{column}{0.48\linewidth}
\begin{block}{Ce que je formaliserais davantage}
\begin{tightitems}
  \item Recrutement avec les RH.
  \item Référentiels de performance individuels.
  \item Plans de formation suivis dans le temps.
  \item Perspectives de carrière et entretiens professionnels.
\end{tightitems}
\end{block}
\end{column}
\end{columns}
\end{frame}
